\scanpage{124}
% archive foliation: 123
\begin{center}
U N I V E R S I T E\\
DES SCIENCES ET TECHNIQUES\\
DU LANGUEDOC\\
\textasciicircum\textasciicircum\textasciicircum\textasciicircum\textasciicircum\textasciicircum\textasciicircum\textasciicircum\textasciicircum\textasciicircum\textasciicircum\textasciicircum
\end{center}
Montpellier, le 4 novembre 1980

\uline{Service de la Scolarité}

AM/gs

Diffusion : Liste A\\
N° 5791

\begin{center}
\uline{AVIS DE SOUTENANCE D'UNE THESE DE DOCTEUR-INGENIEUR}

\uline{``SCIENCES DE L'EAU''}
\end{center}

\uline{NOM et prénom du candidat} : PAPPALARDO Alain

\uline{Titre de la thèse} : ``Contribution à l'étude in situ des transferts de masse dans la zone non saturée des sols''

\uline{Date, heure et lieu de soutenance prévus} : Le 12 novembre 1980 à 10 heures

Salle de cours n° 9 - Université des Sciences et Techniques du Languedoc

\uline{Nom du Directeur de thèse} : M. C. BOCQUILLON

\uline{Désignation du Laboratoire où la thèse est préparée} : Laboratoire d'Hydrologie Mathématique. U.S.T.L. Montpellier.

\uline{Analyse de l'ouvrage} : Cette étude des phénomènes dispersifs lors d'un transfert de soluté dans un milieu poreux se décompose en trois parties:

a) Mise au point du principe d'utilisation d'une sonde prototype gamma/neutrons, susceptible de détecter la présence de deutérium dans l'eau du sol. Etalonnage et test de la méthode lors d'expériences de transferts d'eau deutérée sur la colonne de laboratoire.

b) Utilisation de cette méthode non destructice d'observation en continu de la solution du sol, pour obtenir des profils de concentration, lors de 5 expériences de transfert d'eau deutérée au champ, sur sol nu, en régime hydrodynamique permanent et transitoire et sur sol cultivé couvert de luzerne, sur le site expérimental de la ferme de Cadarache (C.E.A.).

c) Essais de modélisation des résultats expérimentaux obtenus in situ. Test de plusieurs modèles numériques du transfert d'eau et de soluté : modèle monodispersif et d'hétérogénéité. Utilisation de solutions analytiques plus simples à mettre en oeuvre.

\uline{Durée approximative des recherches ayant abouti à la thèse} : 3 ans

\uline{Composition du jury proposé} :\\
MM. Cl. BOCQUILLON\\
Ph. COUCHAT\\
JP. GAUDET\\
E. SERVAT

P/ le Secrétaire Général,\\
l'Attachée d'Administration\\
Universitaire,

\uline{\unsure{Gentil}}

\uline{A. GENTIL}
